\documentclass{article}
\usepackage{geometry}
\usepackage{graphicx} % Required for inserting images
\usepackage{amsmath, amsthm, amssymb}
\usepackage{parskip}
\usepackage{xr}
\usepackage{amsmath}
\usepackage{amssymb}
\newgeometry{vmargin={15mm}, hmargin={24mm,34mm}}
\theoremstyle{definition} 
\newtheorem{definition}{Definition}

\newtheorem{theorem}{Theorem}[section]
\newtheorem{lemma}[theorem]{Lemma}
\newtheorem{corollary}{Corollary}[theorem]
\newtheorem{proposition}[theorem]{Proposition}
\newtheorem{example}[theorem]{Example}

\newcommand{\N}{\mathbb{N}}
\newcommand{\Z}{\mathbb{Z}}
\newcommand{\R}{\mathbb{R}}
\newcommand{\Q}{\mathbb{Q}}
\newcommand{\C}{\mathbb{C}}
\newcommand{\Nil}{\text{Nil}}
\newcommand{\gauss}{\Z[i]}

\title{Galois Theory}
\date{April 2024}
\author{Boran Erol}

\begin{document}

\maketitle

\section{Resources}

\begin{itemize}
    \item The Basic Graduate Year In Algebra
    \item Richard Elman, Lectures on Abstract Algebra, Chapters 11 and 12
\end{itemize}

\section{Definitions and Basic Properties}

\begin{definition}
    A \textbf{Fermat prime} is $f_{n} = 2^{2^{n}} + 1$ such that $f_{n}$ is prime.
\end{definition}

$f_{5}$ is not prime. $f_{1},f_{2},f_{3},f_{4}$ are prime. It's unknown whether we have
infinitely many non Fermat primes or Fermat primes. It's an extremely difficult number
theoretical problem.

If $n = 2^{n}p_{1} \ldots p_{k}$ where $p_{i}$ is a Fermat prime, we can construct an $n$-gon.

Here's something that's often useful in field theory: every field $F$ has a prime subfield
that's generated by $1$. $F$ is a vector space over it's prime subfield, and this consideration
often leads to useful results.

\subsection{Applications of Galois Theory}

Langlands Program: Galois groups of extensions of $\Q$ have "something to do" with modular forms.
Proving bits of "something to do" leads to striking results. This is also used in the proof
of Fermat's Last Theorem.

The Galois group turns out to be related to the fundamental group in algebraic topology.
Field extensions end up corresponding to covering spaces. Algebraic closures end
up corresponding to universal covering spaces.

Inverse Problem: Given a finite group $G$, is there a Galois extension $K$ of $\Q$ such that Galois Group
of $K$ over $\Q$ is $G$? This is an open problem though it's been proven for specific types of groups such
as solvable and Abelian groups.

\begin{lemma}\label{bezout_trick}
    Let $F$ be a field, $n,m$ be coprime and assume $\alpha^{n}, \alpha^{m} \in F$.
    Then, $\alpha \in F$.
\end{lemma}
\begin{proof}
    By Bezout's Identity, $\exists r,s \in \Z: rn + sm = 1$.
    Then, $\alpha = {\alpha^{n}}^{r} + {\alpha^{m}}^{s}$, so
    $\alpha \in F$.
\end{proof}

\newpage

\section{Splitting Fields}

Though the splitting field is unique upto isomorphism, the choice of isomorphism is not unique.
Therefore, there's a minor inconvenience when talking about splitting fields.

\begin{lemma}\label{x_to_p_minus_a_irreducible_iff_no_roots}
    Let $p$ be a prime, $F$ be a field and $a \in F$.
    Then, $f(x) = x^{p} - c$ is irreducible in $F[x]$
    if and only if $x^{p} - c$ has no root in $F$.
\end{lemma}
\begin{proof}
    Assume $F$ has characteristic $p$. Let $L$ be the splitting field of $f$
    and $\alpha \in L$ such that $f(\alpha) = 0$. 
    
    Then, $\alpha^{p} = c$ and $f = x^{p} - \alpha^{p} = {(x - \alpha)}^{p} \in L[x]$.

    If $\alpha \in F$, $f(x) = {(x - \alpha)}^{p}$ in $F[x]$ as well, which is clearly not irreducible.

    Now, let's prove that $\alpha \notin F \implies f$ is irreducible in $F[x]$.

    Suppose not. Then, $f = gh$ for some $g,h \in F[x]$ with $g$ irreducible.
    Then, $g = (x - \alpha)^{m}$ in $L[x]$, so the constant coefficient of $g$
    is $\alpha^{m}$, which has to lie in $F$. Similarly, $\alpha^{p} \in F$.
    Then, by Lemma \ref{bezout_trick}, $\alpha \in F$, which is a contradiction.

    Now, assume $F$ has characteristic $0$. 
\end{proof}

\newpage

\section{Classical Straightedge and Compass Constructions}


Construction by ruler and compass. You can't measure.

The ruler just allows you to draw straight lines between two points on the plane.

The compass allows you to draw a circle with some radius $r$.

Euclid

Here are some problems they considered:
\begin{enumerate}
    \item Constructing a sqaure with the area of a circle
    \item Trisecting an angle
    \item Constructing a regular polygon
\end{enumerate}

In general, it's not true that if an extension has degree that's equal to a power of two,
that it's a square root tower.

\textbf{Squaring the Circle}

Given $z_{1} = 0$, $z_{2} = 1$, we'd like to construct a line segment of length $\sqrt{\pi}$.
Since $\pi$ is transcendetal, this is impossible.

\newpage

\section{Finite Fields}

\newpage

\section{Frobenius Endomorphism}

\begin{lemma}
    Let $F$ be a field of characteristic $p$. Prove that $x \mapsto x^{p}$ is an injective map.
\end{lemma}
\begin{proof}
    Assume $x,y$ map to the same value. Then, $(x-y)^{p} = x^{p} - y^{p} = 0$.
    Since $F$ is a domain, this implies that $x - y = 0$.
\end{proof}



\newpage

\section{Normal Extensions}

Let $K/F$ be algebraic.

Suppose $p \in F[x]$ is irreducible with $\alpha \in K$
such that $p(\alpha) = 0$. Are all roots of $p$ in $L$?

We have seen that this is false in general, since $\Q(\sqrt[3]{2})$
doesn't contain all roots of $x^{3} - 2$.

There are three equivalent conditions of being a normal extension, which the next theorem
explores.

\begin{theorem}
    The following conditions are equivalent for any extension $K/F$:
    \begin{enumerate}
        \item Any irreducible polynomial $p \in F[x]$ that has a root in $K$ factors into linear factors in $K[x]$.
        \item $K$ is the splitting field over $F$ of some set of polynomials $S \subseteq F[x]$.
        \item Let $L$ be an algebraic closure with $F \subset K \subset L$. Any automorphism of $L$ fixes $K$.
    \end{enumerate}
\end{theorem}
\begin{proof}
    Let's first prove that 1 implies 2. This is easy, since we can just associate an irreducible $p_{\alpha} \in F[x]$
    for every $\alpha \in L$. Then, $K$ is the splitting field of the set of these polynomials.

    Let's now prove that 2 implies 3.

    Let $\sigma$ be an automorphism $L$. Then, $\sigma$ fixes all polynomials in $S$ since the coefficients
    of these polynomials are in $F$. Since $L$ is the set of roots of these polynomials, $L$ is also fixed by
    $\sigma$.

    Let's now prove that 3 implies 1. This is the trickiest implication to prove.

    Suppose that any automorphism of $L$ fixes $K$ and let $p \in F[x]$ be an irreducible
    polynomial such that $\alpha \in K$ and $p(\alpha) = 0$.

    We claim that for any root $\beta$ of $p$ in $K$, there's $\sigma(Aut(L))$ such that
    $\sigma(\alpha) = \beta$. This isomorphism exists since $K(\alpha)$ and $K(\beta)$
    are isomorphic and the isomorphism is given by mapping $\alpha$ to $\beta$, as proven in class.
    Since any automorphism $\sigma$ fixes $K$, this implies that $\beta \in K$.
\end{proof}

\begin{definition}
    An extension $K/F$ is a \textbf{normal extension} if it satisfies any of the three
    properties above.
\end{definition}

Normal extensions are easy to construct: pick your favorite irreducible
polynomial and adjoin all of its roots.

\begin{lemma}
    Let $K/F$ be an extension. If $[K:F] = 2$, $K/F$ is a normal extension.
\end{lemma}
\begin{proof}
    
\end{proof}



\newpage

\section{Separable Extensions}

\subsection{Different Resources for Separable Extensions}

\begin{itemize}
    \item Richard Borcherd Lectures
    \item VisualMath -- What is Separable and Normal Extensions?
    \item Elman Algebra Book
    \item Dummit and Foote
    \item Hungerford Algebra
    \item Wikipedia Page
\end{itemize}


\subsection{Definition and Motivation}

Defintiions of separability are in Notability.

\begin{lemma}\label{separable_iff_algebraic_and_min_poly_is_separable}
    $\alpha \in K/F$ is separable if and only if $\alpha$ is algebraic and $m_{\alpha}(t)$ is separable.
\end{lemma}


\begin{lemma}
    Let $f \in F[t]$ and $charF = 0$. $f^{\prime} = 0$ if and only if $f$ is a constant polynomial.
\end{lemma}
\begin{proof}
    The reverse direction is obvious. The forward direction can be proven using the contrapositive.
\end{proof}

\begin{lemma}
    Let $f \in F[t]$ and $charF = p > 0$. Then, $f(t)^{p} = f^{p}(t^{p})$.
    If $f^{\prime} = 0$, then there exists a polynomial $g \in F[t]$ such that $f(t) = g(t^{p})$.
\end{lemma}
\begin{proof}
    The first claim follows from an inductive application of 
    the equality $(a+b)^{p} = a^{p} + b^{p}$.

    \textbf{Prove the second claim, but it's kind of obvious with contradiction.}
\end{proof}

\begin{lemma}\label{characterization_of_perfect_fields_of_characteristic_p}
    Let $F$ be a field of characteristic $p > 0$ and $\sigma$ be the Frobenius endomorphism.
    $F$ is perfect if and only if $\sigma$ is surjective.
\end{lemma}

The following lemma uses the exact same ideas as the previous lemma.

\begin{lemma}\label{idk_what_to_call_this}
    Let $K/F$ be an algebraic extension. If $\alpha \in K$ is separable over $F(\alpha^{p})$,
    $\alpha \in F(\alpha^{p})$.
\end{lemma}
\begin{proof}
    Consider $f(x) = x^{p} - \alpha^{p} = {(x - \alpha)}^{p} \in F(\alpha^{p})[x]$. Notice that $f(\alpha) = 0$.
    Then, $f$ can't be irreducible in $F(\alpha^{p})[x]$ since otherwise $\alpha$ wouldn't be separable as $f$
    has multiple roots.

    Then, let $f = gh$ for some $g,h \in F(\alpha^{p})$ where $g$ is irreducible. Notice that $g = {(x - \alpha)}^{m}$,
    where $1 \leq m < p$. Notice that the last coefficient of $g$ is $-\alpha^{m}$, which implies that $\alpha^{m} \in F^{\alpha^{p}}$.
    
    Then, $\gcd(p,m) = 1$, so there's some $r,s \in \Z$ such that $1 = rp + sm$ by the Bezout Identity. Then, $\alpha = {(\alpha^{p})}^{r} + 
    {(\alpha^{m})}^{s}$, so $\alpha \in F({\alpha^{p}})$.
\end{proof}

\begin{lemma}
    Every finite field is perfect.
\end{lemma}
\begin{proof}
    By \ref{characterization_of_perfect_fields_of_characteristic_p}, we only need to show that the Frobenius endomorphism $\sigma$
    is surjective. However, any injective automorphism of a finite set is surjective, so we conclude the proof.
\end{proof}

\begin{lemma}
    Let $F$ with $charF = p > 0$ and $K/F$ be a separable field extension.
    Then, $K = F(K^{p})$, where $K^{p} := \{ x^{p}: x \in K\}$.
\end{lemma}
\begin{proof}
    It suffices to show $K \subset F(K^{p})$, the reverse containment is trivial.

    Let $\alpha \in K$. Then, $\alpha$ is separable over $F$, so $\alpha$ is separable over $F(\alpha^{p})$.
    Then, by Lemma \ref{idk_what_to_call_this}, $\alpha \in F(\alpha^{p}) \subset K^{p}$. We thus conclude
    the proof.
\end{proof}

\begin{lemma}
    Let $F$ be a field of characteristic $p > 0$.

    Assume $K/F$ is finite and $K = F(K^{p})$. If $\{x_{1}, \ldots, x_{n}\} \subset K$ is linearly independent
    over $F$, then so is $\{{x_{1}}^{p}, \ldots, {x_{n}}^{p}\}$.
\end{lemma}

\begin{lemma}\label{sufficient_condition_for_separability_over_fields_of_characteristic_p}
    Let $F$ be a field of characteristic $p > 0$. If $K/F$ is finite and $K = F(K^{p})$, then $K/F$ is separable.
\end{lemma}

\begin{lemma}\label{adjoining_separable_elements_produce_separable_extensions}
    Let $K/F$ be an extension of fields. If $\alpha \in K$ is separable over $F$,
    $F(\alpha)/F$ is separable.
\end{lemma}
\begin{proof}
    If $F$ has characteristic $0$, the statement is immediately true
    because every algebraic extension of a field of characteristic $0$ is separable. Thus,
    assume without loss of generality that $F$ has characteristic $p > 0$.

    Since $\alpha$ is separable over $F$, $\alpha$ is also separable over $F(\alpha^{p})$.
    Then, by Lemma \ref{idk_what_to_call_this}, $\alpha \in F(\alpha^{p})$. Then, $F(\alpha) = F(\alpha^{p})$.

    Then, by Lemma \ref{sufficient_condition_for_separability_over_fields_of_characteristic_p}, $F(\alpha)/F$
    is separable.
\end{proof}

\begin{lemma}
    If $\alpha_{1}, \ldots, \alpha_{n} \in K$ are separable over $F$, then, $F(\alpha_{1}, \ldots, \alpha_{n})/F$
    is separable.
\end{lemma}
\begin{proof}
    Notice that if an element $x \in K$ is separable over $F$, it's separable over any extension
    of $F$ as well. In particular, $\alpha_{i+1}$ is separable over $F(\alpha_{1},\alpha_{2},\ldots,\alpha_{i})$.
    
    Therefore, $F(\alpha_{1},\alpha_{2},\ldots,\alpha_{i},\alpha_{i+1})/F(\alpha_{1},\alpha_{2},\ldots,\alpha_{i})$
    is a separable extension by Lemma \ref{adjoining_separable_elements_produce_separable_extensions}.

    A simple inductive argument using $K/F$ is separable if and only if $K/E$ and $E/F$ is separable concludes the proof.
\end{proof}

\begin{lemma}\label{separable_part_of_an_extension_is_a_field}
    Let $F_{sep} = \{ \alpha \in K: \alpha \text{ separable over $F$} \}$. Then, $F_{sep}$ is a field.
\end{lemma}
\begin{proof}
    Let $\alpha, \beta \in F_{sep}$. It suffices to show that $\alpha + \beta, \alpha \beta, \alpha^{-1} \in F_{sep}$.

    Since $F(\alpha)/F$ is separable and $\alpha^{-1} \in F(\alpha)$, $\alpha^{-1}$ is separable.

    Since $F(\alpha, \beta)/F$ is separable and $\alpha + \beta, \alpha \beta \in F(\alpha, \beta)$, 
    $\alpha + \beta$ and $ \alpha \beta $ are also separable.

    We thus conclude the proof.
\end{proof}

\begin{lemma}
    Let $F_{0}$ be a field of positive characteristic $p$. Let $F = F_{0}(t_{1}^{p}, t_{2}^{p})$ and 
    $L = F_{0}(t_{1}, t_{2})$. Then, if $\theta \in L \backslash F$, $[F(\theta): F] = p$.
\end{lemma}
\begin{proof}
    Let $\theta \in L \backslash F$. Then, $\theta^{p} \in F$.
    Then, $\theta$ satisfies $f = x^{p} - \theta^{p} = (x - \theta)^{p}$.
    $f$ is irreducible since otherwise $\theta \in F$ using the Bezout trick.
    Thus, $[F(\theta) : F] = p$.
\end{proof}

\subsection{Hungerford's Treatment}

\begin{definition}\label{purely_inseparable_extensions}
    Let $K/F$ be a field extension. An algebraic element $\alpha \in K$ is \textbf{purely inseparable over F}
    if its irreducible polynomial $m_{\alpha, F}(x) \in F[x]$ factors in $K[x]$ as $f = (x - \alpha)^{m}$.

    $K$ is a \textbf{purely inseparable extension of $K$} if every element of $K$ is purely inseparable
    over $F$.
\end{definition}

\begin{lemma}\label{there_are_no_nontrivial_separable_and_inseparable_elements}
    Let $K/F$ be an extension. $\alpha \in K$ is both separable and purely inseparable
    if and only if $\alpha \in F$.
\end{lemma}
\begin{proof}
    $\alpha \in K$ is separable and purely inseparable over $F$ if and only if 
    its irreducible polynomial, $(x - \alpha)^{m}$, has $m$ distinct roots
    in some splitting field. Clearly, this happens if and only if $m = 1$.
\end{proof}

If $F$ has characteristic zero, every algebraic element over $F$ is separable over $F$.
Therefore, the only purely inseparable elements are elements of $F$. Thus, purely inseparable
extensions are non-trivial if $charF > 0$.

\begin{lemma}\label{alpha_separable_implies_all_powers_are_separable}
    Let $K/F$ be an extension with $charF = p > 0$. If $\alpha \in K$
    is separable over $F$, $\alpha^{n}$ is separable over $F$
    for any $n \in \N$.
\end{lemma}
\begin{proof}
    Suppose not. Then, $m_{\alpha^{n}, F}(t)$ is not separable.
    Notice that $m_{\alpha^{n},F}(t^{n})$ has $\alpha$ as a root,
    and is not separable, which is a contradiction.
\end{proof}

\begin{lemma}\label{powers_of_elements_are_separable_in_fields_of_char_p}
    Let $K/F$ be an extension with $charF = p > 0$.
    If $\alpha \in K$ is algebraic over $F$, $\alpha^{p^{n}}$ is separable over $F$ for some $n \geq 0$.
\end{lemma}
\begin{proof}
    We can assume without loss of generality that $\alpha$ is not separable
    using Lemma \ref{alpha_separable_implies_all_powers_are_separable}.
    We induct on the degree of $\alpha$, which we denote by $n$. The base case
    with $n = 1$ is obvious, since $\alpha \in F$. Suppose $n > 1$.

    Let $f$ be the minimal polynomial of $\alpha^{p^{n}}$.
    Then, $\exists g \in F[t]: f(t) = g(t^{p})$. Notice that
    the degree of $g$ is less than the degree of $f$ and that $\alpha^{p^{n-1}}$
    is a root of $g$. Thus, $g$ is separable. Then, $f$ is separable.
\end{proof}

\begin{theorem}\label{characterizations_of_purely_inseparable_extensions}
    Let $K/F$ be an extension with $charF = p > 0$.
    Then, the following statements are equivalent:

    \begin{enumerate}
        \item $K$ is purely inseparable over $F$.
        \item The irreducible polynomial of any $\alpha \in K$ is of the form $x^{p^{r}} - a \in K[x]$.
        \item If $\alpha \in K$, then $\alpha^{p^{n}} \in F$ for some $n \geq 0$.
        \item $K_{sep} = F$.
        \item $K$ is generated over $F$ by a set of purely inseparable elements.
    \end{enumerate}
\end{theorem}
\begin{proof}
    Notice that 1 and 5 are equivalent trivially.

    Using Lemma \ref{there_are_no_nontrivial_separable_and_inseparable_elements}, 4 and 5 are
    also equivalent.

    It remains to show that the conditions 1,2,3 are equivalent.

    Let's first prove that 1 implies 2.

    Assume $K$ is purely inseparable over $F$.
    Let $\alpha \in K$ and $f = {(x - \alpha)}^{m}$ be the irreducible
    polynomial of $\alpha$. Let $m = n p ^{r}$ with $\gcd(n,p) = 1$.
    
    Then, 
    
    \[{(x - \alpha)}^{m} = {(x - \alpha)}^{n p ^{r}} = {(x^{p^{r}} - \alpha^{p^{r}})}^{n}\] 

    Then, it suffices to show that $n = 1$. To prove that $n = 1$,
    it suffices to show that $\alpha^{p^{r}} \in F$. As $f$
    is irreducible, this will imply that $n = 1$. Notice that
    the coefficient of $x^{p^{r}(n-1)}$ in $f$ is $\pm n \alpha^{p^{r}}$,
    which is in $K$ by definition. Since $\gcd(p,n) = 1$, we have that
    $\alpha^{p^{r}} \in F$ and we thus conclude the proof.

    Notice that 2 implies 3 is trivial.

    Let's now prove that 3 implies 1.
    Let $\alpha \in K$. We'd like to show that the minimal polynomial $f$ of $\alpha$ in $F[x]$
    has the form ${(x - \alpha)}^m$ for some $m$. Since $\alpha^{p^{n}} \in F$ for some $n > 0$,
    $\alpha$ is a root of $x^{p^{n}} - \alpha^{p^{n}} = (x - \alpha)^{p^[n]}$. Since $f$
    divides this polynomial, we're done with the proof.
\end{proof}

\begin{lemma}
    Let $K/F$ be an algebraic field extension. Then, $K/F_{sep}$ is purely inseparable
    and $K_{psep} := \{ \alpha \in K: \alpha \text{ is purely inseparable over $F$}\}$
    is a field.
\end{lemma}
\begin{proof}
    If $charF = 0$, $F_{sep} = K$ and thus $K/K$ is a trivial extension.

    Thus, assume $charF = p > 0$. Then, by Lemma \ref{powers_of_elements_are_separable_in_fields_of_char_p},
    $\exists n > 0: \alpha^{p^{n}} \in F_{sep}$. Then, the third condition from Lemma \ref{characterizations_of_purely_inseparable_extensions}
    finishes the proof.
\end{proof}

\newpage

\section{Galois Extensions}


\end{document}
